\begin{abstract}
Curtailment of variable renewable generation is zero-inflated, heavy-tailed and drifts as the system changes, a demanding target for calibrated uncertainty. We study it in Chile's National Electricity System, where about 19{,}000~GWh were curtailed between 2022 and mid-2026. We release an open, errata-corrected, plant-level dataset of daily and hourly curtailment (53 months, 300 plants) that reconciles exactly with the operator's annual totals, and use it to evaluate conformal calibration of seven-day-ahead upper limits under a documented distributional change. With a lognormal hurdle base model, an adaptive scheme combining score transport and online recalibration narrows the limits of the most divergent window by a third at near-nominal coverage. Across nine base models from two families the gain does not transfer, and it is largely mechanical: a constant-threshold oracle that matches the adaptive coverage without adapting in time explains \cf{mec_r2_pct}\% of the variation in width change across 45 model-window cells, as a controlled simulation confirms. The gain measures the drift of the base model's calibration, which here depends more on the base model than on the window. Improving the base model dominates: early stopping alone lowers the hurdle's interval score by \cf{es_vs_ref_static_IS_pct}\%. On high-curtailment days defined by the outcome every method loses coverage, and the transport-based scheme, which leaves more curtailed energy uncovered than a static split, is not the method we recommend for operation. A static conformal wrapper around a multi-quantile predictive gives the lowest operational cost among methods with near-nominal coverage. Dataset and code are public.
\end{abstract}
